\scanpage{22}
\deleted{\unsure{par analog}} \inserted{des} permutations paires sur l'ens des sommets (ou, dualement, sur l'ens des côtés). \deleted{Donc} On en conclut que plus gén. $G_s$ opère \inserted{\unsure{fidèlt|fidèlem$^{\mathrm t}$|fidèlement}} sur $T$ par permutations paires. Aussi, on trouve que \deleted{pour} tt polygône combinatoire de rang $\deleted{\unsure{rg|n}}\equiv1\ (4)$, l'ens des sommets est canoniquement orienté, via épinglage qui permet de numéroter l'ens des sommets (et l'orientation de l'ens des sommets ne dépend pas de l'épinglage\ldots)

On peut \unsure{au|aussi|donc} dire que tt $s\in S$ définit sur $T$, via $P(s)\xrightarrow{\sim}T$, une structure polygonale (pentagonale) \inserted{$\pi_s$}, qui est invariante sous $G_s$ (qui opère exactement comme le groupe des automorphismes de $\pi_s$. A fortiori, $G_s$ invarie l'orientation correspondante de $T$, soit $\omega_s$, donc $G_s$ opère par permutations paires. [On verra plus bas que $\omega_s$ ne dépend pas en fait du choix de $s\in S$]

Fixons-nous un repère $r=\{s,a,f\}$ de $\Pi$, d'où des automorphismes $\sigma_0,\sigma_1,\sigma_2$ engendrant $G=\mathrm{Aut}(\Pi)$.
% archive sheet number: 21
